% =====================================================================
% Relatório Técnico — Atividade 2 (AV2) — Projeto e Análise de Algoritmos
% PROCC/UFS 2026.2
%
% Compilação (Overleaf ou distribuição TeX Live / MiKTeX):
%
%     pdflatex relatorio_av2.tex
%     pdflatex relatorio_av2.tex      % 2x, para numeração de páginas/refs
%
% Requer figuras/tempo_por_n_W1024.png (gráfico experimental da AV2).
% =====================================================================
\documentclass[11pt,a4paper]{article}

\usepackage[T1]{fontenc}
\usepackage[utf8]{inputenc}
\usepackage[brazilian]{babel}
\usepackage{lmodern}
\usepackage[a4paper,top=2.2cm,bottom=2.6cm,left=2.2cm,right=2.2cm]{geometry}
\usepackage{xcolor}
\usepackage{graphicx}
\usepackage{amsmath}
\usepackage{amssymb}
\usepackage{placeins}
\usepackage{booktabs}
\usepackage{array}
\usepackage{tabularx}
\usepackage{enumitem}
\usepackage{listings}
\usepackage{float}
\usepackage{mdframed}
\usepackage{titlesec}
\usepackage{fancyhdr}
\usepackage{lastpage}
\usepackage{amsthm}
\usepackage{algorithm}
\usepackage{algpseudocode}
\usepackage[hidelinks]{hyperref}
% ---------- cores ----------
\definecolor{azul}{HTML}{1F4E79}
\definecolor{azul2}{HTML}{2E5C8A}
\definecolor{laranja}{HTML}{C55A11}
\definecolor{verde}{HTML}{2E7D32}
\definecolor{fundo}{HTML}{F1F4F8}
\definecolor{fundolaranja}{HTML}{FFF6EF}
\definecolor{fundoverde}{HTML}{F0F7F0}
\definecolor{fundoazul}{HTML}{F3F7FC}
\definecolor{linha}{HTML}{C7CED6}
\definecolor{cinza}{HTML}{555555}

% ---------- espaçamento de parágrafo ----------
\setlength{\parindent}{0pt}
\setlength{\parskip}{6pt}
\linespread{1.05}

% ---------- títulos coloridos ----------
\titleformat{\section}{\normalfont\Large\bfseries\color{azul}}{\thesection.}{0.5em}{}
\titleformat{\subsection}{\normalfont\large\bfseries\color{azul2}}{\thesubsection}{0.5em}{}
\titlespacing*{\section}{0pt}{16pt}{6pt}

% ---------- rodapé ----------
\pagestyle{fancy}
\fancyhf{}
\renewcommand{\headrulewidth}{0pt}
\renewcommand{\footrulewidth}{0.4pt}
\fancyfoot[L]{\footnotesize\textcolor{cinza}{Projeto e Análise de Algoritmos --- AV2 $\cdot$ PROCC/UFS 2026.2}}
\fancyfoot[R]{\footnotesize\textcolor{cinza}{Página \thepage\ de \pageref{LastPage}}}

% ---------- caixas ----------
\newmdenv[linecolor=laranja,linewidth=3pt,topline=false,rightline=false,bottomline=false,
  backgroundcolor=fundolaranja,innerleftmargin=10pt,innerrightmargin=10pt,
  innertopmargin=8pt,innerbottommargin=8pt,skipabove=8pt,skipbelow=8pt]{nota}
\newmdenv[linecolor=verde,linewidth=3pt,topline=false,rightline=false,bottomline=false,
  backgroundcolor=fundoverde,innerleftmargin=10pt,innerrightmargin=10pt,
  innertopmargin=8pt,innerbottommargin=8pt,skipabove=8pt,skipbelow=8pt]{teo}
\newmdenv[linecolor=azul,linewidth=1pt,backgroundcolor=fundoazul,
  innerleftmargin=10pt,innerrightmargin=10pt,
  innertopmargin=8pt,innerbottommargin=8pt,skipabove=8pt,skipbelow=8pt]{inv}
\newmdenv[linecolor=laranja,linewidth=1pt,backgroundcolor=fundolaranja,
  innerleftmargin=10pt,innerrightmargin=10pt,
  innertopmargin=8pt,innerbottommargin=8pt,skipabove=8pt,skipbelow=8pt]{status}

% ---------- pseudocódigo (conteúdo em ASCII para compilar sem tropeços) ----------
\lstdefinestyle{pseudo}{
  basicstyle=\ttfamily\footnotesize,
  backgroundcolor=\color{fundo},
  frame=single, framesep=6pt, rulecolor=\color{linha},
  keywordstyle=\color{azul}\bfseries,
  commentstyle=\color{cinza}\itshape,
  morekeywords={funcao,para,ate,faca,se,entao,senao,enquanto,retorne},
  columns=fullflexible, breaklines=true, showstringspaces=false,
  xleftmargin=4pt, xrightmargin=4pt, aboveskip=8pt, belowskip=8pt,
}
\lstset{style=pseudo}
% ---------- teoremas e pseudocódigo da Seção 6 ----------
\theoremstyle{plain}
\newtheorem{lema}{Lema}
\newtheorem{teorema}{Teorema}
\newtheorem{proposicao}{Proposição}
\newtheorem{corolario}{Corolário}
\theoremstyle{definition}
\newtheorem{exemplo}{Exemplo}
\newtheorem{observacao}{Observação}
\theoremstyle{plain}
\renewcommand{\proofname}{Demonstração}
\floatname{algorithm}{Algoritmo}
\algrenewcommand\algorithmicrequire{\textbf{Entrada:}}
\algrenewcommand\algorithmicensure{\textbf{Saída:}}
\algrenewcommand\algorithmicfunction{\textbf{função}}
\algrenewcommand\algorithmicif{\textbf{se}}
\algrenewcommand\algorithmicthen{\textbf{então}}
\algrenewcommand\algorithmicelse{\textbf{senão}}
\algrenewcommand\algorithmicfor{\textbf{para}}
\algrenewcommand\algorithmicdo{\textbf{faça}}
\algrenewcommand\algorithmicreturn{\textbf{devolva}}
\algrenewcommand\algorithmicend{\textbf{fim}}

\newcolumntype{Y}{>{\raggedright\arraybackslash}X}

\hypersetup{pdftitle={Relatório Técnico AV2 - PAA 2026.2},
  pdfauthor={Equipe PAA UFS 2026.2}}

% =====================================================================
\begin{document}

% ---------------------------- CAPA ----------------------------
\begin{titlepage}
\thispagestyle{empty}
\centering
\vspace*{0.6cm}
{\normalsize\color{cinza}
Universidade Federal de Sergipe --- \textbf{\color{azul}PROCC}\\[2pt]
Projeto e Análise de Algoritmos $\cdot$ Período 2026.2\\[2pt]
Módulo II --- Otimização, Grafos e Processamento de Cadeias \par}

\vfill
{\color{laranja}\rule{0.7\textwidth}{2pt}}\par
\vspace{10pt}
{\large\color{laranja}\textbf{Atividade 2 (AV2) --- Relatório Técnico}}\par
\vspace{14pt}
{\Huge\bfseries\color{azul}
Recuperação e seleção otimizada de contexto
para IA generativa em materiais educacionais abertos\par}

\vspace{10pt}
{\large\color{cinza}
Programação dinâmica e mochila 0/1 para seleção
de chunks sob restrição de tokens:\\
um estudo comparativo com estratégia gulosa\par}
\vspace{8pt}
{\color{azul}\rule{0.7\textwidth}{2pt}}\par

\vspace{20pt}
{\large
João Cosme Sena Sá $\cdot$ Eduardo Henrique do Lago Silva $\cdot$ Helena Carvalho Leal\\[4pt]
Hernandison da Silva Bispo $\cdot$ Nadianne Maria dos Santos Galvão $\cdot$ Rafael Takeguma Goto\par}

\vspace{18pt}
\begin{tabular}{@{}r@{\quad}l@{}}
\textbf{Tema}            & 8 --- Materiais educacionais abertos\\
\textbf{Corpus}          & \emph{Pense Python} (tradução livre da 3\textsuperscript{a} ed. de \emph{Think Python})\\
\textbf{Repositório}     & \small\href{https://github.com/joaocss/PAA_UFS_2026_2_SenaSa_Joao_Bispo_Hernandison_Galvao_Nadianne_Goto_Rafael_Leal_Helena_Silva_Eduardo}{Repositório GitHub da equipe}\\

\textbf{Licença do código} & MIT $\cdot$ Commit de referência registrado\\

\textbf{Vídeo da atividade} & \small\href{https://drive.google.com/drive/folders/1qF5mMkBG4x6ZlZMEypKzlrl6P9b71Dy7?usp=sharing}{Pasta no Google Drive (AV2)}\\
\end{tabular}

\vfill
{\footnotesize\color{cinza}
Documento revisado em \today{} $\cdot$ São Cristóvão, SE\\
Entrega dos artefatos até 04/11/2026 $\cdot$ Apresentação em 05/11/2026\par}
\end{titlepage}

% Versão para entrega: relatório sem sumário, conforme limite de páginas.
% \tableofcontents


% ============================ 1 ============================
\section{Identificação da equipe}

A AV2 é desenvolvida por seis discentes do
Programa de Pós-Graduação em Ciência da
Computação (PROCC/UFS), na disciplina
Projeto e Análise de Algoritmos, período
2026.2.

As responsabilidades abrangem especificação
formal (Parte A), implementação e corretude
(Parte B), análise assintótica (Parte C),
experimentação (Parte D) e relação com IA
generativa (Parte E).

A distribuição individual está registrada
em \texttt{docs/CONTRIBUICOES.md} e
resumida na Seção~\ref{sec:contrib},
com distinção entre responsabilidades
acordadas e evidências de execução.
A revisão de código, das provas e dos
resultados é uma atividade coletiva.


% ============================ 2 ============================
\section{Tema, corpus e fonte}

O trabalho aborda o Tema 8 do enunciado:
\textbf{materiais educacionais abertos}.
O corpus é o livro \emph{Pense Python},
tradução livre da terceira edição de
\emph{Think Python}, de Allen B. Downey,
realizada por Rodrigo Castelan Carlson
\cite{downey2024thinkpython}.

O recorte contém 21 notebooks Jupyter,
incluindo prefácio, introdução e capítulos
1 a 19. A aquisição foi registrada em
02/09/2026, utilizando o commit
\texttt{cfde2c450bf4f2ea65326}\allowbreak\texttt{da88e9968f400597e92}.

O material está disponível em
\url{https://rodrigocarlson.github.io/PensePython3ed/},
com licença CC BY-NC-SA 4.0.
O corpus não contém dados sensíveis.

\subsection{Continuidade da Atividade 1}

A AV2 preserva o corpus, a pontuação
TF-IDF e a configuração de recuperação
C3 da AV1, composta por busca linear
e Merge Sort.

A nova etapa seleciona um subconjunto
de chunks sob orçamento de unidades
lexicais, modelado como mochila 0/1.
A programação dinâmica constitui a
abordagem principal, comparada a uma
estratégia gulosa por relevância/custo
\cite{kleinberg2005algorithm}.

São reutilizadas 30 consultas definidas
na AV1, armazenadas em
\texttt{data/queries.csv}, garantindo
a comparação sobre o mesmo conjunto
de perguntas.

O texto integral do corpus não é
redistribuído no repositório.
O script \texttt{scripts/baixar\_corpus.py}
permite reproduzir a aquisição no
commit fixado, enquanto
\texttt{data/corpus\_manifest.csv}
registra os hashes SHA-256 para
verificação de integridade.


% ============================ 3 ============================
\section{Repositório, licença e reprodutibilidade}

O código-fonte está disponível no
repositório público:

\begin{quote}
\footnotesize
\href{https://github.com/joaocss/PAA_UFS_2026_2_SenaSa_Joao_Bispo_Hernandison_Galvao_Nadianne_Goto_Rafael_Leal_Helena_Silva_Eduardo}{Repositório público da equipe (GitHub)}
\end{quote}


O código utiliza licença MIT
(\texttt{LICENSE}), distinta da licença
CC BY-NC-SA 4.0 do corpus.

A versão de referência consultada para
a elaboração deste relatório corresponde
à branch \texttt{av2-parte-e-rag},
no commit:

\begin{quote}
\small\ttfamily
cfe609c21631db6f585b39b51cf8df166d6f07c1
\end{quote}


% ============================ 4 ============================



% ==========================================================
% PARTE A — ESPECIFICAÇÃO DO PROBLEMA
% ==========================================================


\section{Especificação do problema}
\label{sec:especificacao}

\subsection{Problema e representação dos candidatos}

A AV2 investiga a seleção de um subconjunto de
chunks recuperados pela configuração lexical C3
da AV1, buscando maximizar a relevância acumulada
sob um orçamento de unidades lexicais. O problema
é modelado como mochila 0/1, em que cada chunk
é selecionado integralmente ou descartado
\cite{kleinberg2005algorithm}.

O corpus corresponde ao livro \emph{Pense Python},
tradução brasileira da terceira edição de
\emph{Think Python} \cite{downey2024thinkpython}.
O recorte da AV1 contém 21 notebooks Jupyter,
segmentados em 2.737 chunks, dos quais 1.405
são de texto Markdown e 1.332 de código.

Cada candidato $c_i$ possui identificador,
texto, custo lexical $w_i$, relevância $v_i$
e metadados, incluindo \texttt{source\_order}.
O custo é obtido de \texttt{n\_tokens}, que
representa unidades lexicais e não necessariamente
tokens de um modelo de linguagem.

Para uma consulta $q$, os candidatos são definidos
por:

\begin{equation}
C_q=\{c_1,c_2,\ldots,c_n\},
\label{eq:candidatos_av2}
\end{equation}

em que $n=|C_q|$ é a quantidade de candidatos
com relevância positiva, $w_i\in\mathbb{N}^{+}$
e $v_i\in\mathbb{R}_{\geq0}$, com valores finitos.
O orçamento máximo é $W\in\mathbb{N}_0$.

A entrada é representada por:

\begin{equation}
I=(q,C_q,W).
\label{eq:entrada_av2}
\end{equation}

\subsection{Relevância lexical}

A pontuação dos candidatos preserva o mecanismo
TF-IDF da AV1, fundamentado em frequência de
termos e frequência inversa de documento
\cite{manning2008information}.

Para uma consulta $q$ e um chunk $c$, a
relevância é calculada por:

\begin{equation}
\operatorname{score}(q,c)=
\sum_{t\in\operatorname{TermosDistintos}(q)}
qtf_t\,\mathrm{tf}(t,c)\,\mathrm{idf}(t),
\label{eq:score_tfidf}
\end{equation}

em que $qtf_t$ é a frequência do termo na consulta,
$\mathrm{tf}(t,c)$ sua frequência ponderada no chunk
e $\mathrm{idf}(t)$ sua frequência inversa documental.

A AV2 mantém a pontuação original, definindo:

\begin{equation}
v_i=\operatorname{score}(q,c_i).
\label{eq:valor_chunk}
\end{equation}

A implementação da Parte B converte esses escores
em valores inteiros por
$v_i^{(q)}=\operatorname{round}(1000v_i)$.
Consequentemente, a otimalidade computacional
refere-se aos valores quantizados, e não
necessariamente à soma exata dos escores
originais. Essa distinção é considerada na
análise de corretude e nas recorrências.

A pontuação TF-IDF constitui uma estimativa
lexical e não garante, isoladamente,
cobertura semântica ou ausência de redundância.

\subsection{Função objetivo e restrições}

A saída é um subconjunto $S\subseteq C_q$,
acompanhado do custo lexical total e da
relevância acumulada:

\begin{equation}
W_S=\sum_{c_i\in S}w_i,
\qquad
V_S=\sum_{c_i\in S}v_i.
\label{eq:saida_av2}
\end{equation}

Define-se uma variável binária $x_i$,
igual a 1 quando o candidato $c_i$ é
selecionado e 0 caso contrário.

A função objetivo da especificação é:

\begin{equation}
\max V_S=\max\sum_{i=1}^{n}v_i x_i.
\label{eq:funcao_objetivo}
\end{equation}

Sujeita às restrições:

\begin{equation}
\sum_{i=1}^{n}w_i x_i\leq W,
\label{eq:restricao_orcamento}
\end{equation}

\begin{equation}
x_i\in\{0,1\},
\quad i=1,\ldots,n.
\label{eq:restricao_binaria}
\end{equation}

O modelo impede a seleção fracionária e
a repetição de candidatos. A restrição de
capacidade utiliza a mesma unidade lexical
para $w_i$ e $W$.

\subsection{Corretude e critérios de desempate}

A corretude da seleção considera duas
propriedades: viabilidade e otimalidade
em relação à função objetivo adotada.

A viabilidade exige:

\begin{equation}
S\subseteq C_q
\quad\text{e}\quad
\sum_{c_i\in S}w_i\leq W.
\label{eq:corretude_viabilidade}
\end{equation}

Para a formulação original, uma solução
ótima $S^*$ satisfaz:

\begin{equation}
\sum_{c_i\in S^*}v_i
\geq
\sum_{c_i\in S'}v_i,
\quad \forall S'\in\mathcal{F},
\label{eq:corretude_otimalidade}
\end{equation}

em que $\mathcal{F}$ é o conjunto de
subconjuntos viáveis.

Na implementação com programação dinâmica,
a garantia de otimalidade se aplica à soma
dos valores quantizados $v_i^{(q)}$.
Entre soluções de igual valor quantizado,
prioriza-se o menor custo total e,
persistindo o empate, a menor sequência
lexicográfica de \texttt{source\_order}.

Casos como candidatos vazios, orçamento
nulo, itens maiores que $W$, identificadores
duplicados e relevância zero são definidos
nos contratos da AV2 e verificados por testes.

\subsection{Estratégia gulosa e protocolo comparativo}

Como linha de base, utiliza-se uma estratégia
gulosa que ordena os candidatos pela razão
entre relevância e custo lexical:

\begin{equation}
r_i=\frac{v_i}{w_i}.
\label{eq:razao_gulosa}
\end{equation}

Na implementação, a razão utiliza os valores
quantizados $v_i^{(q)}$. Os candidatos são
examinados em ordem decrescente de densidade
e selecionados quando cabem no orçamento
restante. Empates são resolvidos por maior
valor e, posteriormente, menor
\texttt{source\_order}.

Embora produza soluções viáveis, essa
estratégia não garante otimalidade global
para a mochila 0/1.


A comparação utilizou os mesmos candidatos
recuperados pela configuração C3 da AV1,
preservando corpus, consultas, custos,
escores e orçamento entre as estratégias.

Foram avaliadas 30 consultas,
$n \in \{10,30,50\}$ candidatos e
$W \in \{256,512,1024\}$ unidades lexicais.
A grade experimental compreendeu 270
cenários, com cinco estratégias de seleção
e cinco repetições cronometradas por cenário
e estratégia, totalizando 6.750 medições.

As métricas incluíram relevância acumulada,
utilização do orçamento, tempo de seleção,
consumo de memória rastreada e razão de
aproximação em relação ao ótimo quantizado.
Os procedimentos e resultados são detalhados
nas seções de metodologia e resultados
experimentais.


% ============================ 5 ============================
\section{Representação dos dados e estratégia de chunking}

A AV2 preserva o processamento do corpus da AV1,
mantendo a normalização, a segmentação e a
representação lexical dos chunks. A recuperação
utiliza a configuração C3 (busca linear e Merge
Sort), sem alterações nos critérios de pontuação
e ordenação.

\subsection{Normalização e segmentação}

Nas células Markdown, aplica-se normalização
NFKC, conversão para minúsculas, tratamento
da pontuação e compactação de espaços,
preservando os acentos. As células de código
mantêm seu conteúdo e indentação originais,
com representação lexical separada.
Na configuração principal, as stopwords
não são removidas e as saídas de execução
são descartadas.

O chunking respeita a hierarquia de notebooks,
capítulos e seções, sem combinar células
Markdown e código. A configuração principal
utiliza tamanho máximo de 256 unidades
lexicais e sobreposição de 32, produzindo
2.737 chunks. As configurações 128/16 e
512/64 foram avaliadas na AV1, mas não
são adotadas como padrão da AV2.

\subsection{Representação computacional}

Cada candidato é representado pelo par
$(w_i,v_i)$, acompanhado de identificador,
\texttt{source\_order} e metadados de origem.
O valor $v_i$ corresponde à relevância TF-IDF
para a consulta, enquanto $w_i$ é obtido
do campo \texttt{n\_tokens}.

O orçamento $W$ limita a soma dos custos
dos chunks selecionados. A contagem lexical
utilizada como custo é uma aproximação,
não necessariamente equivalente aos tokens
de um modelo de linguagem.

A utilização de um único livro limita a
generalização dos resultados. Além disso,
a relevância lexical pode ser influenciada
por termos técnicos em inglês, e a soma
dos escores não modela diretamente
redundância ou cobertura semântica.

% ============================ 6 ============================


% =========================== 6 ===========================
\section{Algoritmos e pseudocódigo}
\label{sec:algoritmos}

A recuperação lexical C3 da AV1 produz os candidatos ranqueados. Na AV2, a seleção
considera o custo lexical $w_i$ e o valor quantizado
$\hat{v}_i=v_i^{(q)}=\operatorname{round}(1000v_i)$ de cada chunk, sob orçamento $W$.
Comparam-se a programação dinâmica (DP) e o guloso por densidade. Nesta seção, $v_i$ é o
escore original da Equação~\eqref{eq:valor_chunk}, $\hat{v}_i$ é o valor que a
implementação otimiza e $\hat{V}_S=\sum_{i\in S}\hat{v}_i$; $W_S$ e $V_S$ seguem a
Equação~\eqref{eq:saida_av2}.

\subsection{Programação dinâmica: mochila 0/1}
\label{sec:dp}

Depois da validação, \texttt{mochila\_dp()} ordena os candidatos por
\texttt{source\_order} crescente e os renumera de $0$ a $n-1$, de modo que
$o_0<o_1<\dots<o_{n-1}$, em que $o_i$ é o \texttt{source\_order} do candidato $i$. Os
valores de \texttt{source\_order} são distintos, e a numeração não depende da ordem em
que a lista chega. Nesta seção, $i$ se refere sempre a essa numeração, e não ao índice de
$C_q$ na Equação~\eqref{eq:candidatos_av2} nem à posição no ranking.

Cada candidato recebe a chave $a_i=\hat{v}_i(W+1)-w_i$ da
Equação~\eqref{eq:chave_recorrencia_dp}, e um subconjunto $S$ recebe
\[
  A(S)=\sum_{i\in S}a_i=\hat{V}_S\,(W+1)-W_S .
\]
Seja $F(i,w)$ a família dos subconjuntos $S\subseteq\{i,\dots,n-1\}$ com $W_S\le w$, que
sempre contém $\emptyset$, e seja $E(i,w)=\max\{A(S):S\in F(i,w)\}$: a maior chave obtida
com o sufixo $i,\dots,n-1$ sem ultrapassar $w$. A tabela é preenchida pela recorrência da
Equação~\eqref{eq:recorrencia_mochila_av2}, com base $E(n,w)=0$. O uso de sufixos permite
que a reconstrução percorra os candidatos do menor para o maior \texttt{source\_order}.

\begin{algorithm}[htbp]
\caption{Mochila 0/1 com os desempates da especificação (\texttt{mochila\_dp})}
\label{alg:dp}
\begin{algorithmic}[1]
\Require candidatos $C_q$ e orçamento $W\in\mathbb{N}_0$
\Ensure seleção $T$, número de células e $\mathit{valor\_otimo}$
\Function{MochilaDP}{$C_q, W$}
  \State \Call{Validar}{$C_q, W$}; $L\gets C_q$ ordenada por \texttt{source\_order}; $n\gets|L|$; $p\gets W+1$
  \State $E[0..n][0..W]\gets 0$ \Comment{a linha $n$ é a base}
  \For{$i\gets n-1$ \textbf{até} $0$}
    \State $a\gets L[i].\hat{v}\cdot p-L[i].w$
    \For{$w\gets 0$ \textbf{até} $W$}
      \State $E[i][w]\gets E[i+1][w]$ \Comment{sem o item $i$}\label{lin:sem}
      \If{$L[i].w\le w$ \textbf{e} $a+E[i+1][w-L[i].w]>E[i][w]$}\label{lin:compara}
        \State $E[i][w]\gets a+E[i+1][w-L[i].w]$ \Comment{com o item $i$}\label{lin:com}
      \EndIf
    \EndFor
  \EndFor
  \State $T\gets\emptyset$;\ $w\gets W$\label{lin:rec-ini}
  \For{$i\gets 0$ \textbf{até} $n-1$} \Comment{do menor ao maior \texttt{source\_order}}
    \State $a\gets L[i].\hat{v}\cdot p-L[i].w$
    \If{$L[i].w\le w$ \textbf{e} $E[i][w]=a+E[i+1][w-L[i].w]$}\label{lin:leva}
      \State $T\gets T\cup\{L[i]\}$;\ $w\gets w-L[i].w$\label{lin:rec-fim}
    \EndIf
  \EndFor
  \State \Return $T$ na ordem do ranking, $(n+1)(W+1)$, $\lfloor (E[0][W]+W)/p\rfloor$\label{lin:retorno}
\EndFunction
\end{algorithmic}
\end{algorithm}

A comparação da linha~\ref{lin:compara} é estrita: no empate, a célula mantém o ramo
``sem o item'', o que não muda o valor de $E[i][w]$. A reconstrução não consulta qual ramo
foi gravado. Ela testa a igualdade da linha~\ref{lin:leva} e, portanto, leva o item sempre
que levá-lo ainda permite alcançar a chave ótima, inclusive no empate. A ordenação final
pela posição no ranking (linha~\ref{lin:retorno}) segue o contrato 11 de
\texttt{docs/CONTRATOS\_AV2.md} e não altera o conjunto nem as somas.

\subsection{Estratégia gulosa por relevância/custo}
\label{sec:guloso}

A linha de base (S3, \texttt{guloso\_densidade}) ordena os candidatos pela razão
$r_i=\hat{v}_i/w_i$, versão quantizada da Equação~\eqref{eq:razao_gulosa}, comparada como fração exata
(\texttt{Fraction}), e percorre todos eles, levando cada um que ainda cabe.

\begin{algorithm}[htbp]
\caption{Guloso por relevância/custo (\texttt{guloso\_densidade}, S3)}
\label{alg:guloso}
\begin{algorithmic}[1]
\Require candidatos $C_q$ e orçamento $W\in\mathbb{N}_0$
\Ensure seleção viável $S$
\Function{GulosoDensidade}{$C_q, W$}
  \State \Call{Validar}{$C_q, W$}
  \State $L\gets C_q$ ordenada por $r_i$ decrescente; empate: maior $\hat{v}_i$, depois menor $o_i$
  \State $S\gets\emptyset$;\ $u\gets 0$ \Comment{$u$: unidades lexicais já usadas}
  \For{cada candidato $c$ de $L$, na ordem}
    \If{$u+c.w\le W$}
      \State $S\gets S\cup\{c\}$;\ $u\gets u+c.w$ \Comment{o que não cabe é ignorado}
    \EndIf
  \EndFor
  \State \Return $S$ na ordem do ranking
\EndFunction
\end{algorithmic}
\end{algorithm}

O módulo implementa também o guloso modificado (S4, \texttt{guloso\_modificado},
contrato 12), que devolve a melhor entre a seleção de S3 e o item isolado de maior valor
que cabe no orçamento. Ele não faz parte da comparação exigida (DP e guloso por
densidade), mas tem garantia de
aproximação (Teorema~\ref{teo:s4}); o pseudocódigo está no Apêndice~\ref{ap:provas}.

\subsection{Justificativa de corretude}
\label{sec:corretude}

A corretude é demonstrada para o valor quantizado $\hat{v}_i$, que é a grandeza otimizada
pela implementação (Seção~\ref{sec:especificacao}). As demonstrações longas estão no
Apêndice~\ref{ap:provas}.

\begin{lema}[Chave]
\label{lem:chave}
Para subconjuntos viáveis $S$ e $S'$, $A(S)>A(S')$ se e somente se
$\hat{V}_S>\hat{V}_{S'}$, ou $\hat{V}_S=\hat{V}_{S'}$ e $W_S<W_{S'}$. Além disso,
$A(S)\ge 0$, com igualdade somente para $S=\emptyset$.
\end{lema}

\begin{proof}
$A(S)-A(S')=(\hat{V}_S-\hat{V}_{S'})(W+1)-(W_S-W_{S'})$ e $|W_S-W_{S'}|\le W$. Como os
valores são inteiros, $\hat{V}_S>\hat{V}_{S'}$ implica $A(S)-A(S')\ge (W+1)-W=1$; o caso
$\hat{V}_S<\hat{V}_{S'}$ é simétrico; e, se $\hat{V}_S=\hat{V}_{S'}$, a diferença é
$W_{S'}-W_S$. Por fim, $A(\emptyset)=0$ e, se $S\ne\emptyset$, $\hat{V}_S\ge 1$ e
$A(S)\ge (W+1)-W_S\ge 1$.
\end{proof}

O lema depende de os valores serem inteiros (com valores reais, uma diferença de valor
menor que $1/(W+1)$ seria superada por uma diferença de custo) e de o custo não passar de
$W$. Por isso a chave só compara subconjuntos viáveis, que são os únicos que a tabela
considera.

\begin{teorema}[Tabela]
\label{teo:tabela}
Ao fim do preenchimento do Algoritmo~\ref{alg:dp}, $E[i][w]=E(i,w)$ para todo
$0\le i\le n$ e $0\le w\le W$.
\end{teorema}

\begin{proof}
Indução em $i$, de $n$ até $0$. Para $i=n$, $F(n,w)=\{\emptyset\}$ e a linha $n$ é
inicializada com zeros. Para $i<n$, a linha $i+1$ já está completa quando a linha $i$ é
calculada. Os subconjuntos de $F(i,w)$ sem $i$ formam $F(i+1,w)$, com chave máxima
$E(i+1,w)$; os que contêm $i$ existem somente se $w_i\le w$ e são $R\cup\{i\}$ com
$R\in F(i+1,w-w_i)$, com chave máxima $a_i+E(i+1,w-w_i)$, pela aditividade da chave. Pela
hipótese de indução, são os valores lidos nas linhas~\ref{lin:sem} a~\ref{lin:com}, que
gravam o maior deles.
\end{proof}

\begin{teorema}[Reconstrução]
\label{teo:reconstrucao}
O conjunto $T$ devolvido pelo Algoritmo~\ref{alg:dp} é viável, tem chave $E(0,W)$ e,
entre as soluções de chave $E(0,W)$, tem a sequência crescente de \texttt{source\_order}
lexicograficamente menor. Pelo Lema~\ref{lem:chave}, $T$ é a solução definida na
Seção~\ref{sec:especificacao}: maior $\hat{V}_S$; no empate, menor $W_S$; persistindo o
empate, menor sequência de \texttt{source\_order}.
\end{teorema}

A demonstração (Apêndice~\ref{ap:provas}) usa o invariante
$A(T_i)+E(i,w^{(i)})=E(0,W)$, em que $T_i$ e $w^{(i)}$ são o conjunto escolhido e a
capacidade restante antes do passo $i$. Se uma solução ótima $U$ diferisse de $T$ pela
primeira vez num índice $j\in U\setminus T$, a linha~\ref{lin:leva} teria levado $j$.

O valor ótimo sai da chave por divisão inteira:
$\lfloor (E(0,W)+W)/(W+1)\rfloor=\hat{V}_T$, porque
$E(0,W)+W=\hat{V}_T(W+1)+(W-W_T)$ e $0\le W-W_T<W+1$; pelo Lema~\ref{lem:chave} e pelo
Teorema~\ref{teo:reconstrucao}, $\hat{V}_T$ é o maior $\hat{V}_S$ viável. Como
$E(0,W)\ge 0$, o numerador nunca é negativo, e o operador \texttt{//} do Python dá esse
quociente. A parcela $W$ é necessária: sem
ela, o quociente é $\hat{V}_T-1$ sempre que $W_T\ge 1$ (mutação 21 do contrato 14).

\begin{corolario}[Quantização]
\label{cor:quantizacao}
Seja $\mathrm{OPT}$ o maior $V_S$ (escores originais) entre as seleções viáveis e
$k_{\max}\le\min\{n,W\}$ o maior número de itens de uma seleção viável. Então
$V_T\ge \mathrm{OPT}-k_{\max}\cdot 0{,}001$.
\end{corolario}

A cota soma o erro de arredondamento, no máximo $0{,}0005$ por item, de duas seleções: a de
$T$ e a ótima nos escores originais (Apêndice~\ref{ap:provas}).

\begin{proposicao}[Guloso]
\label{prop:guloso}
A seleção do Algoritmo~\ref{alg:guloso} é viável, mas não necessariamente ótima.
\end{proposicao}

\begin{proof}
O laço mantém o invariante $u=W_S\le W$: vale no início, e um item só entra quando
$u+c.w\le W$. Para a não otimalidade, com $W=50$ e
$(w_i,\hat{v}_i)=(10,60),(20,100),(30,120)$ \cite{cormen2022introduction}, o guloso leva os
dois primeiros e obtém $160$, enquanto o segundo e o terceiro somam $220$.
\end{proof}

\begin{teorema}[Guloso modificado]
\label{teo:s4}
Seja $\widehat{\mathrm{OPT}}$ o maior $\hat{V}_S$ entre as seleções viáveis. A seleção de
S4 satisfaz $\hat{V}_{S4}\ge\widehat{\mathrm{OPT}}/2$, e a constante $1/2$ não pode ser
melhorada.
\end{teorema}

Já o S3 sozinho não tem garantia constante: com $W\ge 2$ e os itens $(1,2)$ e $(W,2W-1)$, ele vale
$2$ e o ótimo vale $2W-1$. As demonstrações estão no Apêndice~\ref{ap:provas}, e os custos
de tempo e memória, na Seção~7.

% ====================== MODELO RAM ======================
\section{Análise no modelo RAM}

Adota-se o modelo RAM, com custo constante
para operações elementares sobre valores
de tamanho limitado. Sejam $n$ o número
de candidatos, $W$ o orçamento lexical
e $k\leq n$ a quantidade de itens selecionados.

A programação dinâmica constrói uma tabela
com $(n+1)(W+1)$ posições. A validação
possui custo esperado $O(n)$ e as ordenações
possuem limite $O(n\log n)$. Assim, no
pior caso:

\begin{equation}
T_{\mathrm{DP}}(n,W)
=\Theta(nW+n\log n),
\qquad
M_{\mathrm{DP}}(n,W)=\Theta(nW).
\end{equation}

O custo é pseudopolinomial, pois depende
do valor numérico de $W$, e não apenas
da quantidade de bits de sua representação.

O guloso por densidade ordena os candidatos
e realiza uma varredura linear. Seu custo
temporal e espacial é limitado por:

\begin{equation}
T_{\mathrm{G}}(n)=O(n\log n),
\qquad
M_{\mathrm{G}}(n)=\Theta(n).
\end{equation}

No pipeline completo, a recuperação lexical
C3 e o Merge Sort acrescentam o custo de
busca e ordenação dos resultados. Sejam
$N$ o número de chunks do corpus,
$P$ os resultados positivos, $Q$ o tamanho
da consulta e $L$ a soma dos comprimentos
dos textos. Os limites superiores são:

\begin{equation}
\begin{aligned}
T_{\mathrm{total,DP}}=O(&NQ+L+
P\log(P+1)\\
&+nW+n\log n),
\end{aligned}
\end{equation}

\begin{equation}
\begin{aligned}
T_{\mathrm{total,G}}=O(&NQ+L+
P\log(P+1)+n\log n).
\end{aligned}
\end{equation}

A preparação inicial das estatísticas
documentais não está incluída no custo
por consulta.

\section{Melhor, pior e caso médio}

\paragraph{Programação dinâmica.}

O preenchimento da matriz executa
exatamente $n(W+1)$ iterações,
independentemente dos valores e custos
dos candidatos ou da solução selecionada.

Assim, a inicialização e o preenchimento
possuem custo
$\Theta((n+1)(W+1))$ em todos
os casos de entrada válida.

O custo adicional das ordenações
depende da disposição dos elementos.
Sob o modelo RAM adotado, os limites
assintóticos são:

\begin{equation}
\begin{aligned}
T_{\mathrm{DP}}^{\mathrm{melhor}}(n,W)
&=\Theta(nW),\\
T_{\mathrm{DP}}^{\mathrm{pior}}(n,W)
&=\Theta(nW+n\log n).
\end{aligned}
\label{eq:casos_dp_av2}
\end{equation}

\paragraph{Estratégia gulosa.}

O percurso de seleção examina
todos os $n$ candidatos. Entretanto,
as ordenações podem apresentar
comportamentos diferentes conforme
a disposição das chaves.

Como o Timsort pode explorar
sequências favoráveis à ordenação,
os limites são:

\begin{equation}
\begin{aligned}
T_{\mathrm{G}}^{\mathrm{melhor}}(n)
&=\Theta(n),\\
T_{\mathrm{G}}^{\mathrm{pior}}(n)
&=\Theta(n\log n).
\end{aligned}
\label{eq:casos_guloso_av2}
\end{equation}

O melhor caso linear não deve ser
presumido para as consultas reais,
pois a recuperação C3 ordena os
candidatos por relevância, enquanto
o guloso os ordena por densidade
de relevância/custo.

\paragraph{Caso médio.}

Uma análise probabilística rigorosa requer
uma distribuição de entradas, não definida
para este corpus. Assim, não se atribui
expressão assintótica específica ao caso
médio; os tempos observados são analisados
na seção de resultados.


\section{Recorrências}

\subsection{Mochila 0/1}

Na implementação da Parte B,
o valor computacional de cada item
é quantizado pela escala $1000$:

\[
v_i^{(q)}
=
\operatorname{round}(1000v_i),
\]

em que $v_i$ é o escore TF-IDF
original e \texttt{round()} segue
o arredondamento do Python.

A chave utilizada pela programação
dinâmica é:

\begin{equation}
a_i=v_i^{(q)}(W+1)-w_i.
\label{eq:chave_recorrencia_dp}
\end{equation}

Essa representação prioriza a maior
soma dos valores quantizados e,
em caso de igualdade, o menor custo.
O desempate por sequência
lexicográfica de \texttt{source\_order}
ocorre durante a reconstrução.

Seja $E(i,w)$ a maior chave
obtida a partir dos candidatos
de $i$ até $n-1$, com capacidade $w$.

A condição inicial é $E(n,w)=0$.
A recorrência é:


\begin{equation}
\small
E(i,w)=
\begin{cases}
E(i+1,w), & w_i>w,\\[2mm]
\max\{E(i+1,w),\,
a_i+E(i+1,w-w_i)\}, & w_i\leq w.
\end{cases}
\label{eq:recorrencia_mochila_av2}
\end{equation}


A implementação calcula cada
estado uma única vez. Sua recorrência
temporal de preenchimento é:

\begin{equation}
\begin{aligned}
C(i,W)&=C(i-1,W)+\Theta(W+1),\\
C(0,W)&=\Theta(W+1).
\end{aligned}
\label{eq:recorrencia_custo_dp}
\end{equation}

Por expansão:

\begin{equation}
C(n,W)=\Theta((n+1)(W+1)).
\label{eq:solucao_recorrencia_custo_dp}
\end{equation}

A otimalidade da implementação
refere-se aos valores quantizados,
não necessariamente à soma
exata dos escores originais.

\subsection{Merge Sort}

Na configuração C3, o Merge Sort
divide os $P$ resultados positivos
em duas partes aproximadamente iguais.

Para $P\geq2$, sua recorrência é:

\begin{equation}
T(P)=T(\lfloor P/2\rfloor)
+T(\lceil P/2\rceil)+\Theta(P).
\label{eq:recorrencia_merge_av2}
\end{equation}

Com $T(0)=T(1)=\Theta(1)$.

Para $P$ potência de dois:

\[
T(P)=2T(P/2)+\Theta(P).
\]

Pelo Teorema Mestre,
$a=2$, $b=2$ e
$f(P)=\Theta(P)$.

Como $P^{\log_b a}=P$,
aplica-se o segundo caso:

\begin{equation}
T(P)=\Theta(P\log P),\quad P\geq2.
\label{eq:solucao_merge_av2}
\end{equation}

O resultado também se aplica
a tamanhos que não são potências
de dois: a recorrência com piso
e teto produz profundidade
$\Theta(\log P)$ e custo total
$\Theta(P)$ por nível,
até fatores constantes.

Incluindo as entradas vazias
e unitárias, pode-se escrever:

\[
T(P)=\Theta(1+P\log(P+1)).
\]

As recorrências complementam a análise no
modelo RAM; scripts e verificações estão
documentados no repositório.


% =========================== 10 ===========================
\section{Metodologia experimental}
\label{sec:metodologia}

\subsection{Configuração experimental}

Os experimentos utilizaram o corpus \textit{Pense Python},
30 consultas da AV1 e a recuperação C3 (busca linear
e Merge Sort), mantida fixa entre as estratégias.

Foram comparados S1 (prefixo do ranking),
S2 (guloso por valor), S3 (guloso por densidade),
S4 (guloso modificado) e S5 (mochila 0/1 por
programação dinâmica). A comparação principal
considerou S3 e S5.

A grade utilizou $n\in\{10,30,50\}$ candidatos e
$W\in\{256,512,1024\}$ unidades lexicais.
Foram avaliados 270 cenários e cinco estratégias,
com cinco repetições cronometradas após dois
aquecimentos, totalizando 6.750 medições.

Uma avaliação complementar calculou a diferença
de valor e a redundância textual, além de 1.620
medições do fluxo recuperação--conversão--seleção
para S3 e S5, com três repetições por configuração.

\subsection{Métricas e reprodutibilidade}

As métricas de qualidade incluíram valor quantizado,
relevância TF-IDF, custo lexical, utilização do orçamento,
razão de aproximação e diferença absoluta
$\Delta V=V_{\mathrm{DP}}-V_{\mathrm{estrategia}}$,
calculada sobre valores quantizados.
A redundância foi medida pela similaridade de Jaccard
entre conjuntos de palavras, com limiar $0{,}8$.

Foram registrados tempo de seleção, tempo do fluxo
completo e pico de alocações Python rastreadas
pelo \texttt{tracemalloc}. A preparação do corpus
foi tratada separadamente do tempo de consulta.

As execuções ocorreram no Fedora Linux, com
Python 3.12.16. Dados brutos, gráficos, parâmetros,
logs e scripts estão em \texttt{experimentos/av2/}
e \texttt{scripts/av2/}. A equivalência dos arquivos
de chunks foi confirmada por SHA-256.
A suíte de testes e as auditorias das medições
originais e complementares foram aprovadas.

% =========================== 11 ===========================
\section{Resultados e gráficos}
\label{sec:resultados}

A avaliação compreendeu 6.750 medições de seleção
e 1.620 medições complementares do fluxo completo.
Todas as estratégias respeitaram os orçamentos
estabelecidos, e as verificações de consistência
não identificaram violações nas condições auditadas.

A Tabela~\ref{tab:av2_resultados} apresenta a
comparação principal para $n=50$ candidatos
e $W=1024$ unidades lexicais.
\begin{table}[htbp]
\centering
\caption{Medianas de tempo e memória e média da razão de aproximação para $n = 50$ e $W = 1024$}
\label{tab:av2_resultados}
\small
\begin{tabular}{lrr}
\hline
\textbf{Métrica} & \textbf{S3} & \textbf{S5} \\
\hline
Mediana do tempo (ms) & 0,3030 & 7,2771 \\
Mediana do pico rastreado (KiB) & 7,34 & 854,07 \\
Aproximação média (\%) & 97,78 & 100,00 \\
\hline
\end{tabular}
\end{table}
A Figura~\ref{fig:tempo_av2}, no
Apêndice~\ref{ap:graficos}, apresenta a
variação do tempo de seleção com $n$,
para $W=1024$.

Os valores de memória correspondem ao pico de
alocações Python rastreadas pelo
\texttt{tracemalloc}, não à memória residente
total do processo.

Nesse cenário, a programação dinâmica (S5)
obteve o valor ótimo da função objetivo
quantizada. O guloso por densidade (S3)
alcançou, em média, 97,78\% desse valor,
com menor tempo de seleção e menor
consumo de memória rastreada.

A avaliação complementar registrou diferenças
de valor ($\Delta V$) entre as estratégias
e a programação dinâmica. Não foram detectados
pares de chunks com similaridade de Jaccard
igual ou superior a $0{,}8$ nos conjuntos
selecionados avaliados. Esse resultado
se refere à redundância lexical, não
à equivalência semântica.

Nas medições do fluxo completo, a recuperação
C3 representou uma parcela predominante do
tempo de processamento. Assim, as diferenças
de custo entre os seletores foram menos
expressivas no tempo total do que nas medições
isoladas de seleção.

Os resultados completos, por consulta,
quantidade de candidatos e orçamento, estão
disponíveis em \texttt{experimentos/av2/},
incluindo dados brutos, tabelas e gráficos.

% =========================== 12 ===========================
\section{Discussão de escalabilidade}

Os resultados experimentais são compatíveis
com a análise assintótica apresentada na
Seção~7. A programação dinâmica possui custo
temporal $\Theta(nW+n\log n)$ e espacial
$\Theta(nW)$, enquanto o guloso por densidade
apresenta tempo $O(n\log n)$ e espaço
$\Theta(n)$, sob as hipóteses do modelo RAM.

Para $n=50$ e $W=1024$, a programação dinâmica
apresentou mediana do tempo de seleção de
$7{,}2771$ ms, contra $0{,}3030$ ms do guloso.
O pico de alocações rastreadas foi de
$854{,}07$ KiB e $7{,}34$ KiB, respectivamente.

Esse comportamento evidencia o compromisso
entre otimalidade e custo computacional.
A programação dinâmica maximiza o valor
quantizado sob o orçamento, enquanto o guloso
obteve aproximação média de $97{,}78\%$
nesse cenário, com menor custo de seleção.

Entretanto, nas medições do fluxo completo,
o custo da recuperação C3 predominou sobre
o da seleção. Portanto, as diferenças
observadas entre os seletores não devem ser
interpretadas como diferenças proporcionais
na latência total do sistema.

A extrapolação para conjuntos maiores de
candidatos e orçamentos mais elevados exige
novos experimentos, especialmente devido
à dependência pseudopolinomial da programação
dinâmica em relação a $W$.



% =========================== 13 ===========================
\section{Relação com IA generativa e RAG}

A abordagem desenvolvida pode integrar a etapa
de preparação de contexto em sistemas de geração
aumentada por recuperação
(\emph{Retrieval-Augmented Generation} -- RAG)
\cite{lewis2020retrieval}. Após a recuperação
lexical C3, os chunks são selecionados sob
um orçamento de unidades lexicais antes da
composição do prompt destinado a um LLM.

A programação dinâmica maximiza a relevância
quantizada dentro do orçamento. Para $n=50$
e $W=1024$, o guloso por densidade alcançou
$97{,}78\%$ do valor ótimo, em média,
com menor tempo e memória de seleção.
Isso indica um compromisso entre qualidade
da função objetivo e custo computacional.

Entretanto, as unidades lexicais utilizadas
não equivalem necessariamente aos tokens de
um LLM. Além disso, a pontuação TF-IDF não
garante relevância semântica, diversidade
ou cobertura completa das evidências.

A seleção depende dos candidatos recuperados:
fragmentos relevantes ausentes do ranking C3
não podem ser recuperados pela etapa de
otimização. A função objetivo também não
modela explicitamente dependências entre chunks.

Como não foram avaliadas respostas geradas
por LLMs, os resultados não demonstram
melhorias na qualidade das respostas, no
custo de inferência ou na latência de geração.
Essas questões permanecem como possibilidades
de investigação futura.




% =============== LIMITACOES ===============
\section{Limitações e ameaças à validade}

A validade interna depende da manutenção das
mesmas consultas, candidatos, escores e
orçamentos entre as estratégias. Embora
os experimentos tenham utilizado condições
controladas e repetições, os tempos medidos
podem sofrer influência do sistema operacional,
da execução do Python e de outras atividades
do ambiente computacional.

Quanto à validade de construto, a pontuação
TF-IDF representa relevância lexical, não
necessariamente relevância semântica.
O orçamento utiliza unidades lexicais,
que podem diferir da tokenização de um LLM.
Além disso, a função objetivo aditiva não
penaliza diretamente a redundância nem
representa dependências entre chunks.

A validade externa é limitada pelo uso de
um único corpus educacional, 30 consultas,
$n\in\{10,30,50\}$ e
$W\in\{256,512,1024\}$.
Consequentemente, os resultados não podem
ser generalizados automaticamente para
outros corpora, orçamentos ou ambientes.

Por fim, a otimalidade da programação dinâmica
refere-se aos valores TF-IDF quantizados,
não necessariamente aos escores originais.
A ausência de avaliação com respostas
geradas por LLMs impede conclusões sobre
fidelidade factual, qualidade da geração
ou redução efetiva de custos de inferência.


% ====================== CONCLUSAO ======================
\section{Conclusão}

Este trabalho investigou a seleção de chunks
para composição de contexto sob orçamento
lexical, utilizando programação dinâmica
para mochila 0/1 e estratégias gulosas,
a partir da recuperação C3 da AV1.

Os experimentos evidenciaram um compromisso
entre otimalidade e custo computacional.
Para $n=50$ e $W=1024$, o guloso por
densidade alcançou aproximação média de
$97{,}78\%$ do valor ótimo quantizado,
com mediana do tempo de seleção de $0{,}3030$ ms,
frente a $7{,}2771$ ms da programação
dinâmica.

A programação dinâmica oferece uma
referência ótima para a função objetivo
quantizada, mas apresenta custo
pseudopolinomial em relação ao orçamento.
As medições do fluxo completo também
indicaram predominância do custo da
recuperação lexical sobre o da seleção.

Conclui-se que o guloso por densidade
constitui uma alternativa de menor custo
nos cenários avaliados, enquanto a
programação dinâmica permite estabelecer
a referência de otimalidade.
A avaliação de respostas geradas por LLMs,
a utilização de tokenização específica
e os critérios semânticos de seleção
permanecem como trabalhos futuros.


% =========================== 15 ===========================
\section{Declaração de uso de IA generativa}

Ferramentas de IA generativa foram utilizadas
como apoio à revisão textual e à estruturação
metodológica da especificação.

O registro de ferramentas, modelos, finalidades,
prompts e sugestões aceitas ou rejeitadas
é mantido em \texttt{docs/USO\_DE\_IA.md}.

A equipe permanece responsável pela revisão
das formulações matemáticas, validação do
código, demonstrações de corretude e
interpretação dos resultados. Sugestões
geradas por IA não constituem, isoladamente,
evidências científicas.


% =========================== 16 ===========================
\section{Contribuição individual}
\label{sec:contrib}

As responsabilidades individuais foram
distribuídas da seguinte forma:
João (Partes B e D); Eduardo (B, C e D);
Helena (D e E); Hernandison (A e E);
Nadianne (A e C); Rafael (B e E).

As evidências de execução, incluindo
commits, arquivos, revisões e resultados,
são documentadas em
\texttt{docs/CONTRIBUICOES.md}.
A avaliação crítica e a revisão do
trabalho são responsabilidades coletivas.


% =========================== 17 ===========================
\section{Vídeo da atividade}

O vídeo de apresentação da AV2 documenta
a implementação, as verificações de corretude
e os resultados experimentais, com participação
de todos os integrantes.

URL pública:
\url{https://drive.google.com/drive/folders/1qF5mMkBG4x6ZlZMEypKzlrl6P9b71Dy7?usp=sharing}.

A mesma referência consta no \texttt{README.md}
e no \texttt{VIDEO.md}.


\begin{thebibliography}{99}
\bibitem{cormen2022introduction}
CORMEN, T. H.; LEISERSON, C. E.; RIVEST, R. L.; STEIN, C.
\emph{Introduction to Algorithms}. 4. ed. MIT Press, 2022.

\bibitem{kleinberg2005algorithm}
KLEINBERG, J.; TARDOS, É. \emph{Algorithm Design}.
Pearson/Addison-Wesley, 2006.

\bibitem{skiena2020algorithm}
SKIENA, S. S. \emph{The Algorithm Design Manual}. 3. ed. Springer, 2020.

\bibitem{sedgewick2011algorithms}
SEDGEWICK, R.; WAYNE, K. \emph{Algorithms}. 4. ed. Addison-Wesley, 2011.

\bibitem{manning2008information}
MANNING, C. D.; RAGHAVAN, P.; SCHÜTZE, H.
\emph{Introduction to Information Retrieval}.
Cambridge University Press, 2008.

\bibitem{lewis2020retrieval}
LEWIS, P. et al. Retrieval-Augmented Generation for Knowledge-Intensive NLP Tasks.
\emph{Advances in Neural Information Processing Systems}, 2020.
Disponível em: \url{https://arxiv.org/abs/2005.11401}.

\bibitem{downey2024thinkpython}
DOWNEY, A. B. \emph{Think Python}. 3. ed. O'Reilly Media, 2024.
Tradução brasileira: \emph{Pense Python}, por Rodrigo Castelan Carlson.
Disponível em: \url{https://rodrigocarlson.github.io/PensePython3ed/}.

\bibitem{pedregosa2011scikit}
PEDREGOSA, F. et al. Scikit-learn: Machine Learning in Python.
\emph{Journal of Machine Learning Research}, v. 12, p. 2825--2830, 2011.
\end{thebibliography}

% O corpo principal termina antes desta quebra.
% Os apêndices começam em página separada.
\clearpage
\appendix


% ====================== APENDICES ======================
\clearpage
\appendix

\section{Gráficos experimentais complementares}
\label{ap:graficos}

Este apêndice apresenta evidências complementares
da avaliação das cinco estratégias de seleção.
Os gráficos consideram diferentes quantidades
de candidatos e orçamentos lexicais.

\begin{figure}[!htbp]
    \centering
    \includegraphics[width=0.72\linewidth]
    {figuras/tempo_por_n_W1024.png}
    \caption{Tempo de seleção em função de $n$,
    para $W=1024$.}
    \label{fig:tempo_av2}
\end{figure}

\begin{figure}[!htbp]
    \centering
    \includegraphics[width=0.72\linewidth]
    {figuras/memoria_por_W_n50.png}
    \caption{Pico de alocações Python rastreadas
    em função de $W$, para $n=50$.}
    \label{fig:memoria_av2}
\end{figure}

\begin{figure}[!htbp]
    \centering
    \includegraphics[width=0.72\linewidth]
    {figuras/razao_aproximacao_W512.png}
    \caption{Razão de aproximação em relação
    ao ótimo quantizado, para $W=512$.}
    \label{fig:aproximacao_av2}
\end{figure}

\begin{figure}[!htbp]
    \centering
    \includegraphics[width=0.72\linewidth]
    {figuras/delta_valor_W1024.png}
    \caption{Diferença de valor quantizado em
    relação à programação dinâmica, para $W=1024$.}
    \label{fig:delta_av2}
\end{figure}

\begin{figure}[!htbp]
    \centering
    \includegraphics[width=0.72\linewidth]
    {figuras/fluxo_completo_W1024.png}
    \caption{Tempo do fluxo recuperação,
    conversão e seleção, para $W=1024$.}
    \label{fig:fluxo_av2}
\end{figure}

\begin{figure}[!htbp]
    \centering
    \includegraphics[width=0.72\linewidth]
    {figuras/redundancia_W512.png}
    \caption{Avaliação de redundância lexical
    dos chunks selecionados, para $W=512$.}
    \label{fig:redundancia_av2}
\end{figure}


As Figuras~\ref{fig:tempo_av2} a
\ref{fig:redundancia_av2}, apresentadas
no Apêndice~\ref{ap:graficos}, documentam
as análises complementares de tempo,
memória, aproximação, diferença de valor,
fluxo completo e redundância lexical.

\clearpage



\clearpage
\section{Demonstrações de corretude}
\label{ap:provas}

Este apêndice completa a Seção~\ref{sec:corretude}. A notação é a da
Seção~\ref{sec:dp}.

\subsection{Subestrutura ótima e reconstrução}

\begin{lema}[Subestrutura ótima]
\label{lem:subestrutura}
Seja $S^\ast\in F(i,w)$ com $A(S^\ast)=E(i,w)$, para $i<n$.
Se $i\notin S^\ast$, então $A(S^\ast)=E(i+1,w)$.
Se $i\in S^\ast$, então $w_i\le w$ e $S^\ast\setminus\{i\}$ atinge $E(i+1,w-w_i)$.
\end{lema}

\begin{proof}
Se $i\notin S^\ast$, então $S^\ast\in F(i+1,w)\subseteq F(i,w)$, e nenhum elemento de
$F(i+1,w)$ tem chave maior que $A(S^\ast)$; logo $A(S^\ast)=E(i+1,w)$.
Se $i\in S^\ast$, então $w_i\le W_{S^\ast}\le w$, e $R=S^\ast\setminus\{i\}\in F(i+1,w-w_i)$.
Se algum $R'\in F(i+1,w-w_i)$ tivesse $A(R')>A(R)$, então $R'\cup\{i\}\in F(i,w)$ e
$A(R'\cup\{i\})=a_i+A(R')>A(S^\ast)$, o que contradiz a maximalidade de $S^\ast$.
\end{proof}

Seja $\mathcal{O}=\{S\in F(0,W): A(S)=E(0,W)\}$. Pelo Lema~\ref{lem:chave}, $\mathcal{O}$
contém exatamente as soluções viáveis de maior $\hat{V}_S$ que, entre elas, têm o menor
$W_S$; todas têm o mesmo valor e o mesmo custo. Para cada $S$, seja $\sigma(S)$ a
sequência crescente dos \texttt{source\_order} de $S$, e escreva $\sigma\prec\tau$ quando
$\sigma$ é lexicograficamente menor que $\tau$.

\begin{lema}[Invariante da reconstrução]
\label{lem:invariante}
Seja $T_i$ o conjunto escolhido antes de examinar o candidato $i$ e $w^{(i)}$ o valor de $w$
nesse momento. Para todo $0\le i\le n$, vale $T_i\subseteq\{0,\dots,i-1\}$,
$w^{(i)}=W-W_{T_i}\ge 0$ e $A(T_i)+E(i,w^{(i)})=E(0,W)$.
\end{lema}

\begin{proof}
Para $i=0$, $T_0=\emptyset$ e $w^{(0)}=W$. Suponha o invariante antes do passo $i$ e
escreva $w=w^{(i)}$. Se a condição da linha~\ref{lin:leva} vale, então
$T_{i+1}=T_i\cup\{i\}$, $w^{(i+1)}=w-w_i\ge 0$ e
$A(T_{i+1})+E(i+1,w-w_i)=A(T_i)+E(i,w)=E(0,W)$.
Se não vale, $T_{i+1}=T_i$ e $w$ não muda; nesse caso $E(i,w)=E(i+1,w)$, seja porque
$w_i>w$, seja porque o máximo da recorrência é atingido pelo ramo sem o item. Nos dois
casos, $A(T_{i+1})+E(i+1,w)=E(0,W)$.
\end{proof}

\begin{lema}[Regra de inclusão]
\label{lem:regra}
No passo $i$, com $w=w^{(i)}$, um conjunto $T_i\cup R$, com $R\subseteq\{i,\dots,n-1\}$,
pertence a $\mathcal{O}$ se e somente se $R\in F(i,w)$ e $A(R)=E(i,w)$. Além disso, a
condição da linha~\ref{lin:leva} vale se e somente se existe $U\in\mathcal{O}$ com
$U\cap\{0,\dots,i-1\}=T_i$ e $i\in U$.
\end{lema}

\begin{proof}
A primeira parte vem do Lema~\ref{lem:invariante}: $T_i\cup R$ é viável se e somente se
$W_R\le w$, e sua chave $A(T_i)+A(R)$ é igual a $E(0,W)=A(T_i)+E(i,w)$ se e somente se
$A(R)=E(i,w)$. Para a segunda, se a condição vale, tome $R'$ que atinge
$E(i+1,w-w_i)$; então $R=\{i\}\cup R'$ está em $F(i,w)$ e tem chave
$a_i+E(i+1,w-w_i)=E(i,w)$, e, pela primeira parte, $U=T_i\cup R$ serve. Reciprocamente, se tal $U$ existe,
$R=U\setminus T_i$ contém $i$ e atinge $E(i,w)$, e o Lema~\ref{lem:subestrutura} dá
$w_i\le w$ e $E(i,w)=a_i+E(i+1,w-w_i)$.
\end{proof}

\begin{proof}[Demonstração do Teorema~\ref{teo:reconstrucao}]
Pelo Lema~\ref{lem:invariante} com $i=n$, $W_T\le W$ e $A(T)=E(0,W)$; cada índice é
examinado uma vez, de modo que nenhum candidato se repete. Logo $T\in\mathcal{O}$.
Seja $U\in\mathcal{O}$, $U\ne T$, e seja $j$ o menor índice que está em apenas um dos dois
conjuntos; então $U\cap\{0,\dots,j-1\}=T_j$. Se $j\in U\setminus T$, o
Lema~\ref{lem:regra} garante que a linha~\ref{lin:leva} levaria $j$; contradição. Logo
$j\in T\setminus U$. As duas sequências coincidem antes de $o_j$, e $\sigma(T)$ tem $o_j$
na posição seguinte. Se $U$ não tivesse índice maior que $j$, seria $U=T_j$ e
$W_U<W_{T_j}+w_j\le W_T$, contradizendo $W_U=W_T$; logo $\sigma(U)$ tem nessa posição
algum $o_{j'}>o_j$, e $\sigma(T)\prec\sigma(U)$.
\end{proof}

\begin{observacao}
Como as soluções de $\mathcal{O}$ têm o mesmo custo e todo $w_i\ge 1$, nenhuma sequência
é prefixo próprio de outra, e o resultado não depende da convenção para prefixos. Deixar o
item de fora no empate, ou decidir pela posição no ranking em vez do
\texttt{source\_order}, produz outra regra de desempate; as mutações que trocam a
regra de desempate (4, 13, 14 e 16 do contrato 14) são detectadas pelos testes.
\end{observacao}

\subsection{Quantização}

\begin{proof}[Demonstração do Corolário~\ref{cor:quantizacao}]
Como $\hat{v}_i=\operatorname{round}(1000\,v_i)$, vale $|\hat{v}_i/1000-v_i|\le 0{,}0005$, e,
somando, $|\hat{V}_S/1000-V_S|\le 0{,}0005\,|S|$ para toda seleção $S$. Seja $S^\ast$
viável com $V_{S^\ast}=\mathrm{OPT}$. Como $T$ maximiza $\hat{V}_S$,
\[
V_T\ \ge\ \frac{\hat{V}_T}{1000}-0{,}0005\,|T|
\ \ge\ \frac{\hat{V}_{S^\ast}}{1000}-0{,}0005\,|T|
\ \ge\ \mathrm{OPT}-0{,}0005\,(|S^\ast|+|T|),
\]
e $|S^\ast|,|T|\le k_{\max}$.
\end{proof}

O produto $1000\,v_i$ é calculado em ponto flutuante; o erro relativo, da ordem de
$10^{-16}$, é desprezível diante de $0{,}0005$. O símbolo $k_{\max}$ evita confusão com o
$k$ da Seção~7, que é o número de itens efetivamente selecionados.

\subsection{Guloso modificado (S4)}

\begin{algorithm}[htbp]
\caption{Guloso modificado (\texttt{guloso\_modificado}, S4)}
\label{alg:s4}
\begin{algorithmic}[1]
\Require candidatos $C_q$ e orçamento $W\in\mathbb{N}_0$
\Ensure seleção viável
\Function{GulosoModificado}{$C_q, W$}
  \State \Call{Validar}{$C_q, W$}; $G\gets$ \Call{GulosoDensidade}{$C_q, W$}
  \State $K\gets\{c\in C_q : c.w\le W\}$ \Comment{itens que cabem sozinhos}
  \If{$K=\emptyset$} \State \Return $\emptyset$ \EndIf
  \State $m\gets$ item de $K$ com maior $\hat{v}$; empate: menor posição no ranking
  \If{$\hat{v}_m>\hat{V}_G$} \State \Return $\{m\}$ \EndIf
  \State \Return $G$ \Comment{no empate, fica S3}
\EndFunction
\end{algorithmic}
\end{algorithm}

Quando $K\ne\emptyset$, $\hat{V}_{S4}=\max\{\hat{V}_G,\hat{v}_m\}$. Retirar da entrada os
candidatos com $w_i>W$ não altera S3, S4 nem $\widehat{\mathrm{OPT}}$: no
Algoritmo~\ref{alg:guloso}, como $u\ge 0$, um desses candidatos falha no teste
$u+c.w\le W$ em qualquer momento, sem alterar $u$; no Algoritmo~\ref{alg:s4}, eles ficam
fora de $K$; e nenhuma seleção viável os contém. Pode-se então supor que todo
$w_i\le W$. Se não sobra candidato, $\widehat{\mathrm{OPT}}=0$ e S4 devolve $\emptyset$.

\begin{proof}[Demonstração do Teorema~\ref{teo:s4}]
Suponha que todo $w_i\le W$ e que há ao menos um candidato. Seja $G$ a seleção de S3 e
numere os candidatos na ordem em que o Algoritmo~\ref{alg:guloso} os examina, de modo que
$r_1\ge\dots\ge r_n$. Se S3 leva todos, $G$ é ótima e
$\hat{V}_{S4}\ge\hat{V}_G=\widehat{\mathrm{OPT}}$.
Caso contrário, seja $t$ o primeiro candidato recusado e $P=\{1,\dots,t-1\}$. Ao examinar
$t$, vale $u=W_P$ e $W_P+w_t>W$. Seja $S$ viável. Os itens de $S\setminus P$ têm razão no
máximo $r_t$, e os de $P\setminus S$, pelo menos $r_t$. Então
\[
\hat{V}_{S\setminus P}\le r_t\,W_{S\setminus P}\le r_t\,(W-W_{S\cap P})
< r_t\,(W_P+w_t-W_{S\cap P}) = r_t\,W_{P\setminus S}+\hat{v}_t
\le \hat{V}_{P\setminus S}+\hat{v}_t ,
\]
em que a desigualdade estrita usa $r_t>0$. Somando $\hat{V}_{S\cap P}$,
$\hat{V}_S<\hat{V}_P+\hat{v}_t$; com $S$ ótima, $\widehat{\mathrm{OPT}}<\hat{V}_P+\hat{v}_t$.
Como S3 aceita todo $P$, $\hat{V}_G\ge\hat{V}_P$; como $w_t\le W$, $t\in K$ e
$\hat{v}_m\ge\hat{v}_t$. Logo
$\widehat{\mathrm{OPT}}<\hat{V}_G+\hat{v}_m\le 2\max\{\hat{V}_G,\hat{v}_m\}=2\hat{V}_{S4}$.

Para ver que $1/2$ não melhora, tome $M\ge 2$, $W=2M$ e os itens $(1,2)$, $(M,2M-1)$ e $(M,2M-1)$:
S3 leva os dois primeiros ($2M+1$), o maior item isolado vale $2M-1$ e o ótimo é $4M-2$;
a razão $(2M+1)/(4M-2)$ tende a $1/2$.
\end{proof}

A demonstração usa apenas que a ordem examinada é de razão não crescente, e não depende
dos desempates. A garantia vale para o valor quantizado. Para o escore original, o
argumento do Corolário~\ref{cor:quantizacao} dá
$V_{S4}\ge \hat{V}_{S4}/1000-0{,}0005\,k_{\max}\ge(\mathrm{OPT}-0{,}0005\,k_{\max})/2-0{,}0005\,k_{\max}$,
e, portanto, $V_{S4}\ge\mathrm{OPT}/2-k_{\max}\cdot 0{,}001$.
O exemplo $(1,2)$ e $(W,2W-1)$ da Seção~\ref{sec:corretude} corresponde ao teste
\texttt{test\_adversarial\_guloso\_densidade\_}\allowbreak\texttt{pode\_ser\_muito\_pior}.

\end{document}